\documentclass[12pt,a4paper]{report}

% =========================
% Packages
% =========================
\usepackage[utf8]{inputenc}
\usepackage[T1]{fontenc}
\usepackage[french]{babel}
\usepackage{geometry}
\usepackage{graphicx}
\usepackage{fancyhdr}
\usepackage{titlesec}
\usepackage{xcolor}
\usepackage{float}
\usepackage{enumitem}
\usepackage[most]{tcolorbox}
\usepackage{hyperref}

% =========================
% Geometry
% =========================
\geometry{
    left=2.5cm,
    right=2.5cm,
    top=2.5cm,
    bottom=2.5cm
}

% =========================
% Colors
% =========================
\definecolor{ensablue}{RGB}{0,51,102}

% =========================
% Hyperref setup (après les autres packages)
% =========================
\hypersetup{
    colorlinks=true,
    linkcolor=ensablue,
    filecolor=magenta,
    urlcolor=blue
}

% =========================
% Header and footer
% =========================
\pagestyle{fancy}
\fancyhf{}
\fancyhead[L]{\small\textit{SGITU G4 — Coordination des transports}}
\fancyhead[R]{\thepage}
\renewcommand{\headrulewidth}{0.5pt}
\setlength{\headheight}{15pt}

% =========================
% Chapter and section formatting
% =========================
\titleformat{\chapter}[display]
{\normalfont\Large\bfseries\color{ensablue}}
{\chaptertitlename\ \thechapter}{10pt}{\Large}
\titlespacing*{\chapter}{0pt}{-20pt}{20pt}

\titleformat{\section}
{\normalfont\large\bfseries\color{ensablue}}
{\thesection}{1em}{}

\titleformat{\subsection}
{\normalfont\normalsize\bfseries\color{ensablue}}
{\thesubsection}{1em}{}

\begin{document}

% ============================================
% PAGE DE GARDE
% ============================================
\begin{titlepage}
    \begin{center}
        \vspace*{0.5cm}
        \includegraphics[width=0.4\textwidth]{image.png}

        \vspace{1cm}

        {\large \textbf{Université Abdelmalek Essaâdi}}\\
        {\large \textbf{École Nationale des Sciences Appliquées de Tétouan}}\\

        \vspace{0.5cm}
        \rule{\textwidth}{1pt}
        \vspace{0.8cm}

        {\color{ensablue}\textbf{FILIÈRE : GÉNIE INFORMATIQUE 2}}\\[0.3cm]
        {\large Module : Microservices}\\

        \vspace{1.2cm}

        \begin{tcolorbox}[
            colback=gray!10,
            colframe=ensablue,
            width=0.9\textwidth,
            arc=8mm,
            boxrule=0.8pt,
            drop shadow
        ]
        \centering
        \Large\bfseries
        Rapport Technique Final\\
        Microservice ``Coordination des transports'' (G4)
        \end{tcolorbox}

        \vspace{1.5cm}

        \begin{minipage}[t]{0.45\textwidth}
            \textbf{Réalisé par :}\\[0.3cm]
            EL ASSAL DOUAE\\
            EL MAKOUDI YOUNES\\
            MAGNE TSAFACK LYDIVINE MERVEILLE\\
        \end{minipage}
        \hfill
        \begin{minipage}[t]{0.45\textwidth}
            \raggedleft
            \textbf{Encadré par :}\\[0.3cm]
            Pr. BESRI\\
            \vspace{0.5cm}
        \end{minipage}

        \vfill

        \rule{\textwidth}{1pt}\\
        \vspace{0.3cm}
        {\large \textbf{Année Universitaire : 2025--2026}}

    \end{center}
\end{titlepage}

% ============================================
% TABLE DES MATIÈRES
% ============================================
\tableofcontents
\newpage

% ============================================
% CHAPITRE 1
% ============================================
\chapter{Présentation du microservice}

\section{Description du sous-système}
Le microservice \textbf{G4 -- Coordination des Transports} est responsable de la
coordination opérationnelle du réseau SGITU. Il centralise la gestion des entités
métier nécessaires à l'exploitation quotidienne : lignes, trajets, arrêts,
horaires, affectations véhicule/ligne, missions et événements de coordination
(retard, déviation, panne, incident, annulation).

Son rôle est de fournir une vue cohérente de l'état opérationnel et de permettre
le pilotage des missions de transport en temps réel.

\section{Positionnement dans SGITU}
Dans l'architecture microservices SGITU, G4 occupe une position centrale entre
les services de planification/exploitation et les services d'intégration temps réel.

\begin{itemize}
  \item \textbf{Exposition API} : endpoints REST exposés via la Gateway G10.
  \item \textbf{Interaction suivi des véhicules G7} : consommation des positions véhicules pour le suivi opérationnel et la détection de situations anormales.
  \item \textbf{Interaction gestion des incidents G9} : consommation Kafka \texttt{incident.transport.topic} et persistance des \textbf{impacts} (\texttt{mission\_incident\_impacts}), sans confondre avec les événements opérationnels G4.
  \item \textbf{Interaction Billetterie G1} : publication du cycle de vie des missions (Kafka) selon le contrat inter-services.
  \item \textbf{Interaction Notifications} : relais d'alertes vers le service de notification.
\end{itemize}

Ainsi, G4 agit comme couche de \textit{coordination métier} entre les données de
référence, les flux temps réel et les actions opérationnelles.

\section{Fonctionnalités métier implémentées}
Les fonctionnalités implémentées à ce stade sont les suivantes :

\begin{enumerate}
  \item \textbf{Gestion du référentiel transport}
  \begin{itemize}
    \item CRUD des lignes ;
    \item CRUD des trajets et consultation des arrêts par trajet ;
    \item CRUD des arrêts et filtrage par ligne ;
    \item CRUD des horaires.
  \end{itemize}

  \item \textbf{Gestion opérationnelle}
  \begin{itemize}
    \item gestion des affectations véhicule/ligne ;
    \item planification et suivi des missions ;
    \item transitions de mission (en cours, clôture, annulation) .
  \end{itemize}

  \item \textbf{Gestion des événements de coordination}
  \begin{itemize}
    \item enregistrement d'événements métier ;
    \item détection de retard, déviation, panne et incident ;
    \item annulation de mission sur décision opérationnelle.
  \end{itemize}

  \item \textbf{Intégration et supervision}
  \begin{itemize}
    \item publication Kafka des événements de cycle de vie mission ;
    \item endpoints de supervision (\texttt{/api/g4/health}, statut opérateur, logs) ;
    \item documentation OpenAPI/Swagger et collection Postman de tests.
  \end{itemize}
\end{enumerate}

% ============================================
% CHAPITRE 2
% ============================================
\chapter{Conception}

\section{Identification des endpoints REST}
Les endpoints du microservice G4 sont exposés principalement sous le préfixe \texttt{/api/g4/**}.  
Le service de notification est exposé via \texttt{/api/notifications/send}.  
Les endpoints de supervision et d'état opérateur sont également disponibles pour le monitoring du système.

\section{Choix des méthodes HTTP}
Le design REST suit les conventions standards :
\begin{itemize}
    \item \texttt{GET} : consultation des ressources ;
    \item \texttt{POST} : création de ressource ou action métier ;
    \item \texttt{PUT} : mise à jour d'une ressource ;
    \item \texttt{DELETE} : suppression d'une ressource.
\end{itemize}

\section{Formats d'échange}
Le format d'échange principal est \textbf{JSON}.  
Les contrats d'entrée/sortie sont portés par des DTOs avec validations.  
L'API est documentée via \textbf{Swagger/OpenAPI}.

\section{Gestion des erreurs}
Les codes HTTP appliqués sont :
\begin{itemize}
    \item \texttt{200} / \texttt{201} : succès ;
    \item \texttt{400} : requête invalide (validation) ;
    \item \texttt{401} : non authentifié ;
    \item \texttt{403} : interdit (rôle insuffisant) ;
    \item \texttt{404} : ressource introuvable ;
    \item \texttt{409} : conflit métier (ex. deux missions \texttt{EN\_COURS} sur le même véhicule) ;
    \item \texttt{500} : erreur interne.
\end{itemize}

\section{Diagrammes UML (conception avancée)}
Les diagrammes PlantUML du projet sont versionnés dans \texttt{docs/diagrams/} :
\begin{itemize}
    \item \texttt{mission-coordination-events.puml} — relation Mission / CoordinationEvent ;
    \item \texttt{mission-incident-vs-coordination.puml} — séparation impact incident G9 vs événements G4 ;
    \item \texttt{architecture-externe-g3-g7.puml} — jonctions G3, G7, G1, G5, G10, Kafka.
\end{itemize}
Ces modèles intègrent le retour pédagogique du 11 mai 2026 (feedback prof).

% ============================================
% CHAPITRE 3 bis — Architecture (insert before impl if needed - actually after conception)
% ============================================
\chapter{Architecture technique et choix technologiques}

\section{Stack retenue}
\begin{itemize}
    \item \textbf{Spring Boot 3.5} : productivité, écosystème mature, intégration native Actuator ;
    \item \textbf{PostgreSQL} : persistance relationnelle, contraintes d'intégrité, compatibilité JPA ;
    \item \textbf{Apache Kafka 3.7} : découplage asynchrone inter-groupes (G7, G9, G1) ;
    \item \textbf{Docker Compose} : orchestration locale, réseau partagé \texttt{sgitu-network} ;
    \item \textbf{Resilience4j} : circuit breaker sur appels G5 (Chaos Monkey) ;
    \item \textbf{Micrometer + Prometheus + Grafana} : observabilité (pilier livraison finale).
\end{itemize}

\section{Justification}
Le microservice G4 combine REST synchrone (CRUD, supervision) et messaging asynchrone (positions véhicules, incidents). PostgreSQL garantit la cohérence des missions ; Kafka évite le couplage fort avec G7/G9 en charge.

% ============================================
% CHAPITRE 3
% ============================================
\chapter{Implémentation}

\section{Architecture en couches}
Le microservice suit une architecture en couches :
\begin{itemize}
    \item \texttt{controller} : exposition REST ;
    \item \texttt{service} : logique métier ;
    \item \texttt{repository} : accès aux données via JPA ;
    \item \texttt{entity} : modèle de persistance ;
    \item \texttt{dto/mapper} : transformation entrée/sortie ;
    \item \texttt{config/security} : sécurité JWT, configuration Kafka et technique.
\end{itemize}

\section{Technologies utilisées}
\begin{itemize}
    \item Java, Spring Boot ;
    \item Spring Web, Spring Data JPA, Spring Security ;
    \item JWT (authentification locale de développement) ;
    \item Kafka (publication/consommation d'événements) ;
    \item PostgreSQL (runtime) ;
    \item Swagger/OpenAPI ;
    \item Docker, Docker Compose.
\end{itemize}
\section{Organisation du projet}
Le projet suit une structure Maven standard avec une séparation claire des responsabilités :
\begin{itemize}
    \item \texttt{controller} : exposition des endpoints REST ;
    \item \texttt{service} : implémentation de la logique métier ;
    \item \texttt{repository} : accès aux données via Spring Data JPA ;
    \item \texttt{entity} : modèle de persistance ;
    \item \texttt{dto} : contrats d'échange API (request/response) ;
    \item \texttt{mapper} : transformation entre entités et DTO ;
    \item \texttt{integration} : échanges avec les services externes (G1, G5, G7, G9) ;
    \item \texttt{config} / \texttt{security} : configuration technique, JWT et sécurité ;
    \item \texttt{exception} : gestion centralisée des erreurs.
\end{itemize}

% ============================================
% CHAPITRE 4
% ============================================
\chapter{Documentation API (Swagger / OpenAPI)}

\section{Documentation complète}
La documentation Swagger UI couvre les endpoints d'authentification, de gestion métier (CRUD + actions), de supervision et d'intégration.

\section{Requêtes et réponses}
Chaque endpoint dispose de payloads JSON structurés :
\begin{itemize}
    \item Request DTO avec validations ;
    \item Response DTO cohérente avec le domaine.
\end{itemize}

\section{Cas d'utilisation couverts}
\begin{itemize}
    \item Création et gestion de lignes, trajets, arrêts, horaires ;
    \item Affectation opérationnelle véhicule/ligne ;
    \item Création et suivi des missions ;
    \item Détection/déclaration d'événements de coordination ;
    \item Clôture et annulation de mission.
\end{itemize}

\section{Accès OpenAPI}
\begin{itemize}
    \item \texttt{/swagger-ui.html}
    \item \texttt{/v3/api-docs}
\end{itemize}

\begin{figure}[H]
  \centering
  \includegraphics[width=0.92\textwidth]{fig_4_1_swagger_overview.png}
  \caption{Vue d'ensemble Swagger UI montrant les endpoints exposés du microservice G4.}
  \label{fig:fig_4_1}
\end{figure}

\begin{figure}[H]
  \centering
  \includegraphics[width=0.92\textwidth]{fig_4_2_swagger_request_response.png}
  \caption{Exécution d'un endpoint dans Swagger avec payload de requête JSON et réponse HTTP.}
  \label{fig:fig_4_2}
\end{figure}

\begin{figure}[H]
  \centering
  \includegraphics[width=0.92\textwidth]{fig_4_3_openapi_json.png}
  \caption{Extrait du contrat OpenAPI (\texttt{/v3/api-docs}) montrant version, paths et sécurité JWT (\texttt{bearerAuth}).}
  \label{fig:fig_4_3}
\end{figure}

% ============================================
% CHAPITRE 5
% ============================================
\chapter{Conteneurisation (Docker)}

\section{Introduction}
Afin d'assurer la portabilité du microservice et la reproductibilité de son exécution,
la conteneurisation repose sur une image Docker dédiée à G4. Les dépendances
d'infrastructure (PostgreSQL et Kafka) sont orchestrées séparément via Docker Compose.

\section{Dockerfile}
Le Dockerfile est de type \textit{multi-stage}, ce qui signifie qu'il contient deux étapes :
\begin{itemize}
    \item une étape de build (compilation et packaging Maven) ;
    \item une étape runtime (exécution du JAR dans une image plus légère).
\end{itemize}
Cette approche réduit la taille de l'image finale et isole les dépendances de compilation de l'environnement d'exécution.

\begin{verbatim}
# Build
FROM eclipse-temurin:17-jdk-alpine AS build
WORKDIR /app
COPY .mvn .mvn
COPY mvnw pom.xml ./
COPY src src
RUN chmod +x mvnw && ./mvnw -q -DskipTests package

# Run
FROM eclipse-temurin:17-jre-alpine
WORKDIR /app
RUN addgroup -S sgitu && adduser -S g4 -G sgitu
COPY --from=build /app/target/service-coordination-transport-0.0.1-SNAPSHOT.jar app.jar
RUN chown g4:sgitu app.jar
USER g4
EXPOSE 8084
ENTRYPOINT ["sh", "-c", "java $JAVA_OPTS -jar /app/app.jar"]
\end{verbatim}

\section{Image Docker}
L'image Docker du service est générée localement avec succès à partir du Dockerfile multi-stage.

\begin{figure}[H]
  \centering
  \includegraphics[width=0.92\textwidth]{fig_5_1_docker_build_success.png}
  \caption{Build Docker réussi du microservice G4 (image générée et taguée localement).}
  \label{fig:fig_5_1}
\end{figure}
\section{Exécution}
Le service est exécutable via Docker Compose et exposé sur le port \texttt{8084}.  
La validation d'exécution est démontrée par :
\begin{itemize}
    \item le statut \texttt{Running} des conteneurs dans Docker Compose ;
    \item l'accès à l'interface Swagger sur \texttt{http://localhost:8084/swagger-ui.html}.
\end{itemize}

% ============================================
% CHAPITRE 6
% ============================================
\chapter{Orchestration (docker-compose)}

\section{Introduction}
L'orchestration locale est unifiée dans un seul fichier \texttt{docker-compose.yml},
ce qui permet de démarrer rapidement un environnement cohérent de test et d'intégration.

\section{docker-compose.yml (contenu exact)}
\begin{verbatim}
services:
  kafka:
    image: apache/kafka:3.7.0
    container_name: sgitu-kafka
    ports:
      - "9092:9092"
    environment:
      KAFKA_NODE_ID: 1
      KAFKA_PROCESS_ROLES: broker,controller
      KAFKA_CONTROLLER_QUORUM_VOTERS: 1@localhost:9093
      KAFKA_LISTENERS: PLAINTEXT://:9092,CONTROLLER://:9093
      KAFKA_ADVERTISED_LISTENERS: PLAINTEXT://localhost:9092
      KAFKA_LISTENER_SECURITY_PROTOCOL_MAP: CONTROLLER:PLAINTEXT,PLAINTEXT:PLAINTEXT
      KAFKA_CONTROLLER_LISTENER_NAMES: CONTROLLER
      KAFKA_OFFSETS_TOPIC_REPLICATION_FACTOR: 1
      KAFKA_TRANSACTION_STATE_LOG_REPLICATION_FACTOR: 1
      KAFKA_TRANSACTION_STATE_LOG_MIN_ISR: 1
      KAFKA_AUTO_CREATE_TOPICS_ENABLE: "true"
    healthcheck:
      test: ["CMD-SHELL", "bash -c '</dev/tcp/127.0.0.1/9092'"]
      interval: 10s
      timeout: 5s
      retries: 12
    networks:
      - sgitu-net

  postgres:
    image: postgres:16-alpine
    container_name: sgitu-g4-postgres
    environment:
      POSTGRES_DB: sgitu_g4
      POSTGRES_USER: g4
      POSTGRES_PASSWORD: g4
    ports:
      - "5434:5432"
    volumes:
      - sgitu_g4_pgdata:/var/lib/postgresql/data
    healthcheck:
      test: ["CMD-SHELL", "pg_isready -U g4 -d sgitu_g4"]
      interval: 5s
      timeout: 5s
      retries: 10
    networks:
      - sgitu-net

  g4-coordination:
    build: .
    container_name: sgitu-g4-coordination
    ports:
      - "8084:8084"
    environment:
      SPRING_PROFILES_ACTIVE: docker
      SERVER_PORT: "8084"
      JPA_DDL_AUTO: update
      SGITU_KAFKA_ENABLED: "true"
      SPRING_KAFKA_BOOTSTRAP_SERVERS: "kafka:9092"
      SGITU_G9_INCIDENT_TOPIC: "${SGITU_G9_INCIDENT_TOPIC:-incident.transport.topic}"
      SGITU_G7_POSITIONS_TOPIC: "${SGITU_G7_POSITIONS_TOPIC:-vehicule-positions}"
      SGITU_G1_MISSION_TOPIC: "${SGITU_G1_MISSION_TOPIC:-missions-lifecycle}"
      SGITU_G9_CONSUMER_GROUP: "${SGITU_G9_CONSUMER_GROUP:-g4-coordination-g9}"
      SGITU_G7_CONSUMER_GROUP: "${SGITU_G7_CONSUMER_GROUP:-g4-coordination-g7}"
      SGITU_G10_URL: "${SGITU_G10_URL:-http://g10-gateway:8080}"
      SGITU_G5_URL: "${SGITU_G5_URL:-http://g5-notifications:8085}"
      SGITU_G1_URL: "${SGITU_G1_URL:-http://g1-billetterie:8081}"
      SGITU_G7_URL: "${SGITU_G7_URL:-http://g7-tracking:8087}"
      SGITU_G9_URL: "${SGITU_G9_URL:-http://g9-incidents:8089}"
      SGITU_INTEGRATION_G5_NOTIFICATION_PATH: "${SGITU_INTEGRATION_G5_NOTIFICATION_PATH:-/api/notifications/send}"
      SGITU_G1_MISSION_LIFECYCLE_PATH: "${SGITU_G1_MISSION_LIFECYCLE_PATH:-/api/internal/missions/events}"
    depends_on:
      postgres:
        condition: service_healthy
      kafka:
        condition: service_healthy
    networks:
      - sgitu-net

volumes:
  sgitu_g4_pgdata:

networks:
  sgitu-net:
    name: sgitu-network
    driver: bridge
\end{verbatim}

\section{Interaction inter-microservices}
Le compose prépare les variables d'environnement nécessaires à l'intégration avec la Gateway (G10), les notifications (G5), la billetterie (G1), le suivi véhicule (G7) et les incidents (G9).

\section{Simulation distribuée}
L'environnement local simule un contexte distribué réaliste avec service applicatif, base PostgreSQL et broker Kafka.

\begin{figure}[H]
  \centering
  \includegraphics[width=0.92\textwidth]{fig_6_1_compose_running_3_3.png}
  \caption{Environnement Docker Compose G4 démarré avec succès : \texttt{sgitu-kafka}, \texttt{sgitu-g4-postgres} et \texttt{sgitu-g4-coordination} en statut \textit{Running (3/3)}.}
  \label{fig:fig_6_1}
\end{figure}

\begin{figure}[H]
  \centering
  \includegraphics[width=0.92\textwidth]{fig_6_2_docker_desktop.png}
  \caption{Docker Desktop : projet \texttt{service-coordination-transport} avec les conteneurs \texttt{sgitu-kafka} (\texttt{apache/kafka:3.7.0}, port \texttt{9092}), \texttt{sgitu-g4-coordination} (port \texttt{8084}) et \texttt{sgitu-g4-postgres} (\texttt{postgres:16-alpine}, port hôte \texttt{5434}) en cours d'exécution.}
  \label{fig:fig_6_2}
\end{figure}

% ============================================
% CHAPITRE 7
% ============================================
\chapter{Sécurité (JWT, Spring Security)}

\section{Authentification (JWT)}
Endpoint de login local : \texttt{POST /api/auth/login}, retour d'un token \texttt{Bearer JWT}.
\section{Protection des endpoints}
Les endpoints publics sont limités à :
\begin{itemize}
    \item \texttt{/api/auth/login} ;
    \item \texttt{/api/g4/health} et \texttt{/api/g4/logs} ;
    \item \texttt{/swagger-ui.html}, \texttt{/v3/api-docs} ;
    \item \texttt{/actuator/health}, \texttt{/actuator/prometheus} (monitoring).
\end{itemize}
Tous les autres endpoints sont protégés par Spring Security + JWT, avec \textbf{re-vérification locale des rôles} (indépendamment de G10).

\section{Gestion des accès (rôles alignés G3)}
Les rôles JWT suivent la convention inter-groupes G3/G10 :
\begin{itemize}
  \item \texttt{ROLE\_OPERATOR} — gestionnaire réseau : CRUD lignes, trajets, arrêts, horaires ;
  \item \texttt{ROLE\_DISPATCHER} — gestionnaire flotte : missions, affectations, événements, notifications ;
  \item \texttt{ROLE\_ADMIN\_G4} / \texttt{ROLE\_ADMIN} — supervision et droits complets.
\end{itemize}

\subsection*{Règles d'accès fonctionnelles}
\begin{itemize}
  \item \textbf{Consultation (\texttt{GET /api/g4/**})} : OPERATOR, DISPATCHER, ADMIN\_G4, ADMIN.
  \item \textbf{Référentiel réseau} (lignes, trajets, arrêts, horaires) : OPERATOR, ADMIN\_G4, ADMIN.
  \item \textbf{Opérationnel flotte} (missions, événements, incident-impacts, notifications) : DISPATCHER, ADMIN\_G4, ADMIN.
  \item \textbf{File notifications PENDING} : ADMIN\_G4, ADMIN.
\end{itemize}

\section{Validation croisée}
Un token JWT émis par G3 (même secret) est accepté par G4 pour prouver l'intégration sécurisée inter-groupes (capture Postman fournie dans \texttt{rapport/captures/}).

\section{Supervision et monitoring}
\begin{itemize}
  \item \texttt{GET /api/g4/health} : DB, Kafka, notifications en attente ;
  \item \texttt{GET /api/v1/operator/status} : état des intégrations G1/G5/G7/G9/G10 ;
  \item \texttt{GET /api/g4/logs} : journal applicatif (public, sans token) ;
  \item \texttt{GET /actuator/prometheus} : métriques pour Prometheus/Grafana.
\end{itemize}

% ============================================
% CHAPITRE 8 — Tests
% ============================================
\chapter{Validation et tests}

\section{Tests automatisés}
La commande \texttt{./mvnw test} exécute les tests JUnit 5 et Spring Boot Test :
\begin{itemize}
    \item \texttt{LigneApiIntegrationTest} — CRUD lignes, health public ;
    \item \texttt{MissionVehicleConflictIntegrationTest} — HTTP 409 conflit véhicule ;
    \item \texttt{SecurityRolesIntegrationTest} — 403 si rôle incorrect (OPERATOR vs DISPATCHER) ;
    \item \texttt{G4ApplicationTests} — chargement du contexte Spring.
\end{itemize}

\section{Collection Postman}
Fichier \texttt{postman/SGITU-G4-Coordination-Transport.postman\_collection.json} :
\begin{itemize}
    \item scénarios métier (lignes, missions, événements) ;
    \item dossier \textbf{99 — Validation croisée} (JWT G3 $\rightarrow$ G4) ;
    \item dossier \textbf{100 — Chaos Monkey} (G5 down, notifications PENDING).
\end{itemize}

Exemples JSON : \texttt{postman/examples/}.

% ============================================
% CHAPITRE 9 — Livraison finale (3 piliers)
% ============================================
\chapter{Livraison finale — trois piliers et résilience}

\section{Pilier 1 — Intégration sans couture}
Stack Docker sur \texttt{sgitu-network} : compose local, \texttt{docker-compose.yml} (profil monitoring, sans G5), intégration dans le monorepo \texttt{SGITU-Microservices/docker-compose.yml}.

\section{Pilier 2 — Observabilité}
Health, logs structurés, Prometheus, dashboard Grafana (\texttt{monitoring/}).

\section{Pilier 3 — Validation croisée}
Preuve Postman : appel G4 avec JWT valide émis par G3 (ou G10).

\section{Chaos Monkey (bonus)}
Appel critique G5 : circuit breaker Resilience4j, réponse \texttt{DEGRADED}, table \texttt{pending\_notifications}, retry automatique. Voir \texttt{docs/SOUTENANCE\_CHAOS\_ORAL.md}.

% ============================================
% CHAPITRE 10 — Conclusion
% ============================================
\chapter{Conclusion}

\section{Bilan}
Le microservice G4 répond aux cas d'usage prévus : référentiel réseau, pilotage flotte, événements de coordination, intégrations Kafka/HTTP, sécurité JWT, conteneurisation Docker Compose, documentation OpenAPI et tests automatisés.

\section{Perspectives}
\begin{itemize}
    \item Déploiement Kubernetes (hors périmètre obligatoire mai 2026) ;
    \item CI/CD GitHub Actions ;
    \item Tests inter-groupes en environnement partagé ENSA.
\end{itemize}

% ============================================
% ANNEXE — État d'avancement (archivé)
% ============================================
\chapter*{Annexe — Historique d'avancement}
\addcontentsline{toc}{chapter}{Annexe — Historique}

\section*{Fonctionnalités finalisées}
\begin{itemize}
    \item Endpoints métiers, persistance, JWT, Kafka, Docker, Swagger, Postman, Resilience4j, monitoring.
\end{itemize}

\section*{Difficultés résolues}
\begin{itemize}
    \item Alignement rôles G3/G4 ; image Kafka ; séparation Incident G9 / Événement G4 ; HTTP 409 conflit véhicule.
\end{itemize}

\end{document}
